\section{Literature Review}
\label{sec:related}

\subsection{Analytical and Optimization-Based Bus Holding Control}
\label{sec:related_holding}

The instability of close-interval transit operations was first formalized by \citet{welding1957instability} and \citet{osuna1972control}, who showed that passenger waiting time scales with $1 + \mathrm{CV}^2(H)$, where $\mathrm{CV}^2(H) = \mathrm{Var}(H)/\mathbb{E}[H]^2$ denotes the squared coefficient of variation of vehicle headways. Early real-time control formulations leveraged emerging AVL data to cast stop-level holding as deterministic or stochastic rolling-horizon quadratic programs minimizing weighted passenger waiting and onboard delay \citep{eberlein2001holding,hickman2001analytic,sun2014models}. However, full corridor-wide mathematical programs with capacity constraints \citep{delgado2012much,sanchez2016realtime} require multi-step arrival predictions across all origin--destination pairs and solve mixed-integer or nonlinear programs at every stop event.

To obtain scalable, closed-form feedback laws, \citet{daganzo2009headway} introduced adaptive forward-headway holding, proving that linear feedback on deviations from a target headway eliminates bunching under mild perturbations. Subsequent extensions incorporated cooperative two-sided spacing using both predecessor and follower positions \citep{daganzo2011cooperation,bartholdi2012self}, virtual reference schedules \citep{xuan2011dynamic}, signalized corridor priority \citep{estrada2016bus}, overtaking and distributed boarding dynamics \citep{wu2017modelling}, predictive headway estimation \citep{andres2017predictive}, and dynamic bus substitution \citep{petit2018substitution}. In a comprehensive comparative evaluation, \citet{berrebi2018comparing} demonstrated that bilateral holding rules achieve strong headway regularity without requiring long-horizon predictions, establishing a widely adopted benchmark in transit operations \citep{cats2011impacts,cats2012guidelines}. Nevertheless, as highlighted in the synthesis of \citet{gkiotsalitis2021atstop}, purely headway-equalizing rules remain blind to onboard passenger occupancy, time-varying scheduled headway transitions, and terminal turnaround slack under finite fleets.

\subsection{Dynamic Timetabling and Fixed-Fleet Corridor Circulation}
\label{sec:related_dispatch}

On timescales of $15$--$60\text{ min}$, passenger demand along urban radial corridors exhibits pronounced directional asymmetry and peak-to-trough variation that cannot be addressed by stop-level holding around a static headway. \citet{newell1971dispatching} proved that when a fixed number of daily vehicle dispatches is allocated across a non-stationary arrival intensity $\Lambda(t)$ to minimize total passenger waiting time without capacity saturation, optimal dispatch frequency scales proportionally to $\sqrt{\Lambda(t)}$ (equivalently, target headway $H^\star(t) \propto \Lambda(t)^{-1/2}$). Modern automated fare collection and APC systems enable rich empirical characterization of these spatio-temporal demand profiles \citep{pelletier2011smart,furth2006service}.

Translating time-varying dispatch rates into executable corridor operations is constrained by vehicle circulation across both directions \citep{ibarra2015planning,ceder2016public}. \citet{gkiotsalitis2019robust} formulated robust timetable optimization subject to fleet size and regulatory layover bounds, showing that aggressive peak headway compression erodes terminal recovery time and amplifies trip-to-trip delay propagation. When real-time holding is deployed concurrently with dynamic dispatch on a fixed-pool corridor, the two authorities become tightly coupled through terminal inventory: every second of en-route holding directly consumes terminal layover slack prior to the vehicle's next opposite-direction trip. Existing optimization frameworks rarely coordinate real-time stop holding with dynamic quota-conserving headway allocation in closed loop without resorting to brittle centralized trajectory forecasts.

\subsection{Reinforcement Learning and Multi-Timescale Control in Transit}
\label{sec:related_rl}

Deep reinforcement learning has emerged as an attractive model-free alternative for real-time bus control under complex stochastic dwell and traffic dynamics. \citet{chen2016marl} and \citet{alesiani2018rl} pioneered multi-agent Q-learning and policy-gradient holding controllers on single-loop corridors. \citet{wang2020dynamic} introduced a multi-agent actor--critic framework with graph-structured vehicle interactions, later extending it to asynchronous event-driven execution \citep{wang2021asynchronous} and distributional risk-sensitive critics \citep{wang2023robust}. \citet{rodriguez2023cooperative} combined holding and stop-skipping in a cooperative DRL policy, while \citet{yu2024hierarchical} explored hierarchical multi-agent RL for bunching mitigation. Most recently, \citet{xu2025systematic} conducted a systematic benchmark of RL bus bunching configurations, cautioning that reward functions dominated by headway equalization frequently inflate in-vehicle travel times and fail to outperform simple analytical baselines on net passenger journey time when evaluated under realistic fleet constraints.

In parallel, the machine learning literature has established principles for temporal abstraction \citep{sutton1999between,nachum2018data}, off-policy maximum-entropy actor--critic learning \citep{haarnoja2018soft,christodoulou2019discrete,chen2021redq}, constrained MDPs via Lagrangian duality \citep{altman1999constrained,miryoosefi2019constraints,stooke2020responsive}, and multi-frequency control \citep{lee2020multiple,metelli2020persistence,holt2025evocontrol}. However, directly transplanting generic hierarchical RL to public transit fails for two domain-specific reasons: (i) temporal action abstraction alone does not prevent high-frequency APC sampling noise and localized demand spikes from corrupting upper-level timetable decisions, and (ii) unconstrained lower-level exploration systematically violates terminal turnaround feasibility on fixed-pool corridors.

\begin{table}[t]
\centering
\caption{Positioning of \texttt{FreqDuet} relative to representative transit control paradigms.}
\label{tab:literature_positioning}
\resizebox{\textwidth}{!}{%
\begin{tabular}{@{}lcccccc@{}}
\toprule
Study / Paradigm & Dynamic & Stop-Event & Demand Band & Fixed-Pool & Onboard-Load & Causal APC/AVL \\
 & Headway & Holding & Separation & Turnaround Cap & Journey Obj. & Identification \\
\midrule
\citet{daganzo2009headway} & -- & Forward & -- & -- & -- & \checkmark \\
\citet{berrebi2018comparing} & -- & Bilateral & -- & -- & -- & \checkmark \\
\citet{sanchez2016realtime} & -- & Rolling QP & -- & Soft / Open & \checkmark & Forecast-dep. \\
\citet{gkiotsalitis2019robust} & Offline / Per. & -- & -- & \checkmark & \checkmark & -- \\
\citet{wang2020dynamic,wang2023robust} & -- & MARL & -- & Ring / Uncap. & Partial & Raw state \\
\citet{yu2024hierarchical} & Regional HRL & MARL & -- & Uncapped & Headway-dom. & Raw state \\
\textbf{\texttt{FreqDuet} (Ours)} & \textbf{Slow SAC} & \textbf{CMDP SAC} & \textbf{LF$\to$U, HF$\to$L} & \textbf{Hard FIFO Mask} & \textbf{Exact Pax-s} & \textbf{4-Way Causal} \\
\bottomrule
\end{tabular}%
}
\end{table}

Table~\ref{tab:literature_positioning} summarizes how FreqDuet bridges these research streams. Unlike single-level holding controllers \citep{daganzo2009headway,berrebi2018comparing,wang2020dynamic} or unconstrained hierarchical RL \citep{yu2024hierarchical}, FreqDuet aligns online frequency-decomposed APC signals with the physical authority of each control layer while enforcing exact daily trip quotas and hard FIFO turnaround budgets under estimator-free passenger journey objectives.
